\chapter{de Rham Cohomology}\label{de-rham}
\chapterstatus{complete}

\section{Introduction}\label{de-rham:section-introduction}

Every exact form is closed, because $d \circ d = 0$. The converse is a question about the shape of the manifold. On $\RR^n$, and on any convex open set, every closed form of
positive degree is exact (the Poincar\'e lemma). On the punctured plane $\RR^2 \setminus \{0\}$ the angle form $(x\,dy - y\,dx)/(x^2 + y^2)$ is closed but not exact: its integral
around the origin is $2\pi$, while the integral of an exact form around a closed curve is $0$. The obstruction comes from the hole at the origin. \emph{De Rham cohomology},
closed forms modulo exact forms, measures such obstructions in every degree:
\[
H^k_{\mathrm{dR}}(M) = \frac{\{\text{closed } k\text{-forms}\}}{\{\text{exact } k\text{-forms}\}} .
\]
The de Rham theorem (Theorem~\ref{de-rham:theorem-de-rham}) says that this analytic invariant is the singular cohomology of $M$ with real coefficients, a topological invariant.
This is the bridge between analysis and topology on which Hodge theory is built: the Hodge theorem selects a canonical harmonic representative in each de Rham class
(Chapter~\ref{hodge-theorem}), and on a complex manifold the type decomposition of harmonic forms gives the Hodge decomposition (Chapter~\ref{hodge-decomposition}).

The chapter goes as follows.
\begin{enumerate}
\item The de Rham complex, first examples, homotopy invariance with an explicit homotopy operator, and the Poincar\'e lemma with an explicit formula for primitives
(Section~\ref{de-rham:section-complex}).
\item The Mayer--Vietoris sequence, with its connecting homomorphism made explicit, and the computation of the cohomology of spheres, punctured planes and projective spaces.
\item The de Rham theorem, periods of $1$-forms, the K\"unneth formula, Poincar\'e duality, and the cohomology of compact Lie groups
(Section~\ref{de-rham:section-de-rham-theorem}).
\item Currents, the forms with distribution coefficients, which contain the submanifolds and are the natural home of cycle classes (Section~\ref{de-rham:section-currents}).
\end{enumerate}
References: \cite{bott-tu}, \cite{warner-foundations}, \cite[Chapter 0]{griffiths-harris}.

\noindent\textbf{Prerequisites.} Chapter~\ref{differential-forms} (Differential Forms and Integration), Chapter~\ref{sheaf-cohomology}
(Sheaf Cohomology), Chapter~\ref{poincare-duality} (Topological Manifolds and Poincar\'e Duality) and Chapter~\ref{distributions}
(Distributions and Sobolev Spaces).

\section{The de Rham complex}\label{de-rham:section-complex}

\begin{definition}\label{de-rham:definition-de-rham}
The \emph{de Rham complex} of a smooth manifold $M$ is $0 \to \Omega^0(M) \xrightarrow{d} \Omega^1(M) \xrightarrow{d} \cdots$, a complex by Theorem~\ref{differential-forms:theorem-exterior-derivative}. Its
cohomology $H^k_{\mathrm{dR}}(M) = \ker(d \colon \Omega^k \to \Omega^{k+1})/\operatorname{im}(d \colon \Omega^{k-1} \to \Omega^k)$ is the \emph{de Rham cohomology}; its elements are classes $[\omega]$ of
\emph{closed} forms modulo \emph{exact} ones. The wedge product makes $H^\bullet_{\mathrm{dR}}(M)$ a graded-commutative $\RR$-algebra, and a smooth map $f \colon N \to M$ induces an algebra homomorphism
$f^* \colon H^\bullet_{\mathrm{dR}}(M) \to H^\bullet_{\mathrm{dR}}(N)$. Replacing $\Omega^\bullet(M)$ by the compactly supported forms $\Omega^\bullet_c(M)$ gives $H^\bullet_{\mathrm{dR},c}(M)$, functorial for
open inclusions (extension by zero) and for proper maps.
\end{definition}

\begin{proposition}\label{de-rham:proposition-h0}
$H^0_{\mathrm{dR}}(M) \cong \RR^{\pi_0(M)}$, the locally constant functions. The wedge product is well defined on classes, and $f^*$ commutes with $d$, hence descends to cohomology.
\end{proposition}

\begin{proof}
A function with $df = 0$ has vanishing differential, hence is constant on connected components (Theorem~\ref{real-analysis:theorem-mean-value} in charts, and connectedness). If $\alpha$, $\beta$ are closed then
$d(\alpha \wedge \beta) = 0$, and if $\alpha = d\gamma$ then $\alpha \wedge \beta = d(\gamma \wedge \beta)$, because $d(\gamma \wedge \beta) = d\gamma \wedge \beta \pm \gamma \wedge d\beta$ and $d\beta = 0$; naturality of $d$ is
Theorem~\ref{differential-forms:theorem-exterior-derivative}.
\end{proof}

\begin{example}\label{de-rham:example-first-computations}
\begin{enumerate}
\item \emph{The real line.} Every $1$-form on $\RR$ is $f\,dx$, and it is exact: $f\,dx = dF$ with $F(x) = \int_0^xf$. So $H^0_{\mathrm{dR}}(\RR) = \RR$ and $H^k_{\mathrm{dR}}(\RR) = 0$ for
$k \geq 1$ (there are no forms of degree $\geq 2$).
\item \emph{The circle.} On $S^1 = \RR/2\pi\ZZ$ every $1$-form is closed, and is $f(\theta)\,d\theta$ with $f$ periodic. It is exact if and only if $\int_0^{2\pi}f = 0$: then
$F(\theta) = \int_0^\theta f$ is periodic and $dF = f\,d\theta$; conversely $\int_{S^1}dF = 0$. So $\int_{S^1} \colon H^1_{\mathrm{dR}}(S^1) \to \RR$ is an isomorphism, and $[d\theta]$ is a basis.
Here $d\theta$ is a closed $1$-form which is not the differential of a function, since $\theta$ is not a function on $S^1$.
\item \emph{The punctured plane.} The angle form $\alpha = (x\,dy - y\,dx)/(x^2 + y^2)$ is closed and not exact on $\RR^2 \setminus \{0\}$
(Counterexample~\ref{differential-forms:counterexample-closed-not-exact}); so $H^1_{\mathrm{dR}}(\RR^2 \setminus \{0\}) \neq 0$. It is the pullback of $d\theta$ by
$(x, y) \mapsto (x, y)/r \in S^1$, and we will see that it spans $H^1$ (Example~\ref{de-rham:example-angle-form}).
\end{enumerate}
\end{example}

\begin{theorem}[Homotopy invariance and the Poincar\'e lemma]\label{de-rham:theorem-homotopy-invariance}
Let $i_0, i_1 \colon M \to M \times \RR$ be $p \mapsto (p, 0)$ and $p \mapsto (p, 1)$. There is a linear map $K \colon \Omega^k(M \times \RR) \to \Omega^{k-1}(M)$ with
\[
i_1^*\omega - i_0^*\omega = d(K\omega) + K(d\omega).
\]
Consequently smoothly homotopic maps induce the same map on de Rham cohomology, homotopy equivalent manifolds have isomorphic de Rham cohomology, and $H^k_{\mathrm{dR}}(U) = 0$ for $k > 0$ and
$H^0_{\mathrm{dR}}(U) = \RR$ for every convex open $U \subseteq \RR^n$ (the \emph{Poincar\'e lemma}).
\end{theorem}

\begin{proof}
\emph{The decomposition.} Let $t$ be the coordinate of $\RR$. In local coordinates $x$ on $M$ and $t$, a $k$-form on $M \times \RR$ is a sum of terms $f(x, t)\,dx^I$ and
$g(x, t)\,dt \wedge dx^J$; so it can be written uniquely as
\[
\omega = \alpha_t + dt \wedge \beta_t,
\]
where $\alpha_t$ is a $k$-form and $\beta_t$ a $(k-1)$-form on $M$, depending smoothly on the parameter $t$ (their coefficients are $f(\cdot, t)$ and $g(\cdot, t)$); this is
independent of the coordinates on $M$. For instance $i_s^*\omega = \alpha_s$, since $i_s^*dt = 0$. Write $d_M$ for the exterior derivative in the $M$-directions at fixed $t$, and
$\partial_t$ for the derivative of the coefficients in $t$. Then $d = d_M + dt \wedge \partial_t$ on forms without $dt$, and
\[
d\omega = d_M\alpha_t + dt \wedge \partial_t\alpha_t + d(dt \wedge \beta_t) = d_M\alpha_t + dt \wedge (\partial_t\alpha_t - d_M\beta_t),
\]
because $d(dt \wedge \beta_t) = -dt \wedge d\beta_t = -dt \wedge (d_M\beta_t + dt \wedge \partial_t\beta_t) = -dt \wedge d_M\beta_t$.

\emph{The homotopy operator.} Define $K\omega = \int_0^1\beta_t\,dt$, the $(k-1)$-form on $M$ whose coefficients are the integrals over $t \in [0, 1]$ of those of $\beta_t$. By the
formula for $d\omega$, its $dt$-part is $\partial_t\alpha_t - d_M\beta_t$, so $K(d\omega) = \int_0^1(\partial_t\alpha_t - d_M\beta_t)\,dt$. On the other hand
$d(K\omega) = \int_0^1d_M\beta_t\,dt$, by differentiation under the integral sign (Proposition~\ref{real-analysis:proposition-parameter-integrals}). Adding, and using the fundamental
theorem of calculus for the coefficients,
\[
d(K\omega) + K(d\omega) = \int_0^1\partial_t\alpha_t\,dt = \alpha_1 - \alpha_0 = i_1^*\omega - i_0^*\omega .
\]

If $F \colon M \times \RR \to N$ is smooth with $F \circ i_0 = f$ and $F \circ i_1 = g$, then for closed $\omega$ on $N$,
$g^*\omega - f^*\omega = d(K F^*\omega)$, so $f^* = g^*$ on cohomology. Continuous homotopies may be replaced by smooth ones: continuous maps between manifolds are homotopic to smooth maps, and continuously
homotopic smooth maps are smoothly homotopic \cite[Chapter 2]{hirsch-differential-topology}. A convex $U$ is smoothly contractible, by the straight-line homotopy, so its de Rham cohomology is that of a
point.
\end{proof}

The homotopy operator gives explicit primitives. For a star-shaped open set $U \subseteq \RR^n$ (with respect to $0$), let $X = \sum_ix^i\partial_i$ be the position vector
field.

\begin{corollary}[Explicit Poincar\'e lemma]\label{de-rham:corollary-explicit-poincare}
Let $U \subseteq \RR^n$ be star-shaped with respect to $0$ and $\omega$ a closed $k$-form on $U$, $k \geq 1$. Then $\omega = d\eta$ with
\[
\eta_x = \int_0^1t^{k-1}\,(\iota_X\omega)_{tx}\,dt, \qquad\text{that is,}\qquad \eta = \sum_I\Bigl(\int_0^1t^{k-1}\omega_I(tx)\,dt\Bigr)\iota_X(dx^I),
\]
where $\iota_X(dx^{i_1} \wedge \cdots \wedge dx^{i_k}) = \sum_m(-1)^{m-1}x^{i_m}\,dx^{i_1} \wedge \cdots \wedge \widehat{dx^{i_m}} \wedge \cdots \wedge dx^{i_k}$.
\end{corollary}

\begin{proof}
Let $\rho \colon \RR \to [0, 1]$ be smooth and nondecreasing, with $\rho = 0$ on $(-\infty, 0]$ and $\rho = 1$ on $[1, \infty)$, and apply
Theorem~\ref{de-rham:theorem-homotopy-invariance} to $F(x, t) = \rho(t)x$, a smooth map $U \times \RR \to U$ since $U$ is star-shaped. Since
$F^*dx^i = \rho(t)\,dx^i + \rho^{\prime}(t)x^i\,dt$, the terms of $F^*(\omega_I\,dx^{i_1} \wedge \cdots \wedge dx^{i_k})$ containing $dt$ are
$\omega_I(\rho x)\,\rho^{k-1}\rho^{\prime}\sum_m x^{i_m}\,dx^{i_1} \wedge \cdots \wedge dt \wedge \cdots \wedge dx^{i_k}$, with $dt$ in the $m$-th place; moving $dt$ to the front gives the sign
$(-1)^{m-1}$. So $\beta_t = \rho(t)^{k-1}\rho^{\prime}(t)(\iota_X\omega)_{\rho(t)x}$ in the notation of the proof of the theorem, and the substitution $s = \rho(t)$ gives
$K(F^*\omega) = \int_0^1s^{k-1}(\iota_X\omega)_{sx}\,ds = \eta$. As $F(\cdot, 1) = \id$, $F(\cdot, 0)$ is constant, whose pullback kills forms of positive degree, and $d\omega = 0$, the
homotopy formula gives $\omega = d\eta$.
\end{proof}

\begin{example}\label{de-rham:example-explicit-primitives}
\begin{enumerate}
\item For the closed $1$-form $\omega = 2xy\,dx + x^2\,dy$ on $\RR^2$, $\iota_X\omega = 2x^2y + x^2y = 3x^2y$, so $\eta = \int_0^13t^2\,dt \cdot x^2y = x^2y$, and indeed $d(x^2y) = \omega$.
\item For $\omega = dx \wedge dy$ on $\RR^2$, $\iota_X\omega = x\,dy - y\,dx$ and $\eta = \int_0^1t\,dt\,(x\,dy - y\,dx) = \frac12(x\,dy - y\,dx)$, the primitive used in
Example~\ref{differential-forms:example-stokes-checks}(1) to compute areas.
\item For the volume form $dx \wedge dy \wedge dz$ of $\RR^3$ the formula gives $\frac13(x\,dy \wedge dz + y\,dz \wedge dx + z\,dx \wedge dy)$.
\item The slit plane $\RR^2 \setminus \{(x, 0) : x \leq 0\}$ does not contain $0$, but it is star-shaped with respect to $(1, 0)$: a segment from $(1, 0)$ to a point off the
$x$-axis meets the axis only at $(1, 0)$. So the angle form is exact there; a primitive is the polar angle with values in $(-\pi, \pi)$.
\end{enumerate}
\end{example}

\begin{theorem}[Mayer--Vietoris]\label{de-rham:theorem-mayer-vietoris}
Let $M = U \cup V$ with $U$, $V$ open. There is a long exact sequence
\[
\cdots \to H^k_{\mathrm{dR}}(M) \to H^k_{\mathrm{dR}}(U) \oplus H^k_{\mathrm{dR}}(V) \to H^k_{\mathrm{dR}}(U \cap V) \xrightarrow{\ \delta\ } H^{k+1}_{\mathrm{dR}}(M) \to \cdots
\]
\end{theorem}

\begin{proof}
The sequence of complexes $0 \to \Omega^\bullet(M) \to \Omega^\bullet(U) \oplus \Omega^\bullet(V) \to \Omega^\bullet(U \cap V) \to 0$, with maps $\omega \mapsto (\omega|_U, \omega|_V)$ and
$(\alpha, \beta) \mapsto \alpha|_{U \cap V} - \beta|_{U \cap V}$, is exact. Injectivity: a form vanishing on $U$ and on $V$ vanishes. Exactness in the middle: if $\alpha = \beta$ on $U \cap V$, they
glue to a form on $M$. Surjectivity: for a partition of unity $\{\psi_U, \psi_V\}$ subordinate to $\{U, V\}$ (Theorem~\ref{smooth-manifolds:theorem-partitions-of-unity}) and
$\gamma \in \Omega^\bullet(U \cap V)$, the forms $\psi_V\gamma$ on $U$ and $-\psi_U\gamma$ on $V$, both extended by zero (the support of $\psi_V$ does not meet $U \setminus V$, and similarly for
$\psi_U$), have difference $\gamma$ on $U \cap V$. The maps commute with $d$, so Theorem~\ref{homological-algebra:theorem-long-exact-sequence} gives the long exact sequence. Following the
construction of the connecting homomorphism, for a closed $\gamma$ on $U \cap V$ we lift it to $(\psi_V\gamma, -\psi_U\gamma)$, apply $d$, and observe that $d(\psi_V\gamma)$ and
$-d(\psi_U\gamma)$ agree on $U \cap V$, since $d(\psi_U\gamma + \psi_V\gamma) = d\gamma = 0$. So
\[
\delta[\gamma] = [\omega], \qquad \omega = \begin{cases} d(\psi_V\gamma) = d\psi_V \wedge \gamma & \text{on } U,\\ -d(\psi_U\gamma) = -d\psi_U \wedge \gamma & \text{on } V,\end{cases}
\]
a closed form on $M$ supported in $U \cap V$.
\end{proof}

\begin{example}\label{de-rham:example-spheres}
$H^k_{\mathrm{dR}}(S^n) \cong \RR$ for $k = 0, n$ and $0$ otherwise ($n \geq 1$). Let $U = S^n \setminus \{N\}$ and $V = S^n \setminus \{S\}$ for the north and south poles. Both are
diffeomorphic to $\RR^n$ by stereographic projection, so their cohomology is $\RR$ in degree $0$ and $0$ otherwise; and $U \cap V$ is diffeomorphic to $\RR^n \setminus \{0\}$, which is
homotopy equivalent to $S^{n-1}$ (by $x \mapsto x/|x|$, with the inclusion of the unit sphere as homotopy inverse, and the homotopy $(x, t) \mapsto x/((1 - t) + t|x|)$), so
$H^\bullet_{\mathrm{dR}}(U \cap V) \cong H^\bullet_{\mathrm{dR}}(S^{n-1})$ by Theorem~\ref{de-rham:theorem-homotopy-invariance}. The Mayer--Vietoris sequence splits into pieces:
\begin{enumerate}
\item for $k \geq 1$ the terms $H^k(U) \oplus H^k(V)$ vanish, so $\delta \colon H^k_{\mathrm{dR}}(U \cap V) \to H^{k+1}_{\mathrm{dR}}(S^n)$ is an isomorphism, that is
$H^{k+1}(S^n) \cong H^k(S^{n-1})$;
\item in low degrees, $0 \to H^0(S^n) \to H^0(U) \oplus H^0(V) = \RR^2 \to H^0(U \cap V) \to H^1(S^n) \to 0$.
\end{enumerate}
For $n = 1$, $U \cap V$ has two components, so $H^0(U \cap V) = \RR^2$, the map $\RR^2 \to \RR^2$, $(a, b) \mapsto (a - b, a - b)$, has rank $1$, and $H^1(S^1) \cong \RR$, $H^0(S^1) = \RR$. For
$n \geq 2$, $U \cap V$ is connected, the map $\RR^2 \to \RR$ is surjective, and $H^1(S^n) = 0$, $H^0(S^n) = \RR$. Induction on $n$ with (1) gives the result: $H^n(S^n) \cong H^{n-1}(S^{n-1}) \cong \cdots \cong H^1(S^1) \cong \RR$,
and all other groups vanish. For the circle, the generator of $H^1$ is the form $d\theta$, which is closed but not exact; its integral over $S^1$ is $2\pi$
(Exercise~\ref{differential-forms:exercise-pullback-computation}). In general $\int_M \colon H^n_{\mathrm{dR}}(M) \to \RR$ is well defined for compact oriented $M$ by Stokes' theorem
(Theorem~\ref{differential-forms:theorem-stokes}).

The connecting homomorphism can be seen at work on the circle. Take $U$ and $V$ two open arcs covering $S^1$, whose intersection consists of two arcs $A$ and $B$, and let $\gamma$ be the
locally constant function equal to $1$ on $A$ and $0$ on $B$. Then $\delta[\gamma]$ is represented by $d\psi_V$ on $A$ (and $0$ elsewhere), a $1$-form supported in $A$. Near the end of $A$
where $V$ stops, $\psi_V = 0$, as the support of $\psi_V$ lies in $V$; near the end where $U$ stops, $\psi_U = 0$, so $\psi_V = 1$. Hence $\int_Ad\psi_V = \pm1$, and $\delta[\gamma]$ is a nonzero
multiple of the generator of $H^1(S^1)$: a ``bump'' $1$-form concentrated on one arc.
\end{example}

\section{The de Rham theorem}\label{de-rham:section-de-rham-theorem}

\begin{theorem}[de Rham]\label{de-rham:theorem-de-rham}
For every smooth manifold $M$ there are natural isomorphisms
\[
H^k_{\mathrm{dR}}(M) \cong H^k(M, \RR_M) \cong H^k(M; \RR),
\]
the first to sheaf cohomology of the constant sheaf, the second to singular cohomology with real coefficients. Under these isomorphisms the wedge product corresponds to the cup product, and the
isomorphism is induced by integration: the class of a closed $k$-form $\omega$ corresponds to the singular cochain $\sigma \mapsto \int_{\Delta^k}\sigma^*\omega$ on smooth simplices.
\end{theorem}

\begin{proof}
Let $\underline{\Omega}^k$ be the sheaf $U \mapsto \Omega^k(U)$. The sequence of sheaves
\[
0 \to \RR_M \to \underline{\Omega}^0 \xrightarrow{d} \underline{\Omega}^1 \xrightarrow{d} \cdots
\]
is exact: exactness at $\underline{\Omega}^0$ is Proposition~\ref{de-rham:proposition-h0}, and exactness at $\underline{\Omega}^k$ for $k \geq 1$ holds on stalks, because every point has a basis of convex
coordinate neighbourhoods, on which closed forms are exact by the Poincar\'e lemma (Theorem~\ref{de-rham:theorem-homotopy-invariance}). Each $\underline{\Omega}^k$ is a module over the sheaf of smooth
functions, hence fine, hence acyclic on the paracompact space $M$ (Example~\ref{sheaf-cohomology:example-fine} and
Theorem~\ref{sheaf-cohomology:theorem-soft-acyclic}). So this is an acyclic resolution of $\RR_M$ and
Proposition~\ref{sheaf-cohomology:proposition-hypercohomology}(3) gives $H^k(M, \RR_M) \cong H^k(\Gamma(M, \underline{\Omega}^\bullet)) = H^k_{\mathrm{dR}}(M)$. The second isomorphism is
Theorem~\ref{sheaf-cohomology:theorem-singular-comparison}, since manifolds are locally contractible and paracompact, and it is compatible with products. The description by integration over smooth simplices
is the standard comparison of the two resolutions; see \cite[Chapter 5]{warner-foundations} or \cite[Chapter III]{bredon-sheaf}.
\end{proof}

\begin{corollary}\label{de-rham:corollary-consequences}
For a compact manifold $M$ the spaces $H^k_{\mathrm{dR}}(M)$ are finite-dimensional, they vanish for $k > \dim M$, and $\dim H^k_{\mathrm{dR}}(M) = b_k(M)$, the $k$-th Betti number. For a compact oriented
$n$-manifold, $H^n_{\mathrm{dR}}(M) \cong \RR$ for each connected component, with the isomorphism given by integration.
\end{corollary}

\begin{proof}
Compact manifolds have finitely generated singular homology \cite[Appendix]{hatcher-at}, so Theorem~\ref{de-rham:theorem-de-rham} and the universal coefficient theorem give finite dimension and the
identification with Betti numbers; the vanishing above the dimension holds because there are no nonzero forms of degree $> n$. For the top degree, $H^n(M; \RR) \cong H_0(M; \RR) \cong \RR$ per component by
Poincar\'e duality (Theorem~\ref{poincare-duality:theorem-poincare-duality}), and integration is nonzero on the class of a volume form
(Corollary~\ref{differential-forms:corollary-volume}).
\end{proof}

Explicitly, if $M_1, \ldots, M_c$ are the connected components of the compact oriented manifold $M$ (finitely many, as they are open and cover the compact $M$), then
$\Omega^\bullet(M) = \prod_i\Omega^\bullet(M_i)$, so $H^n_{\mathrm{dR}}(M) = \bigoplus_iH^n_{\mathrm{dR}}(M_i)$, and the isomorphism of the corollary is
\[
H^n_{\mathrm{dR}}(M) \to \RR^c, \qquad [\omega] \mapsto \Bigl(\int_{M_1}\omega, \ldots, \int_{M_c}\omega\Bigr).
\]
So a top degree form on $M$ is exact if and only if its integral over \emph{each} component vanishes; the vanishing of the total integral $\int_M\omega$ is not enough when
$M$ is disconnected. For example, on $M = S^2 \sqcup S^2$ the form equal to the area form on the first sphere and to minus the area form on the second has total integral $0$ but is
not exact.

\begin{proposition}\label{de-rham:proposition-periods}
Let $M$ be a connected manifold and $x_0 \in M$. There is an isomorphism
\[
H^1_{\mathrm{dR}}(M) \xrightarrow{\ \sim\ } \Hom(\pi_1(M, x_0), \RR),
\]
sending $[\omega]$ to the homomorphism whose value on the class of a piecewise smooth loop $\gamma$ is the \emph{period} $\int_\gamma\omega$. In particular $H^1_{\mathrm{dR}}(M) = 0$ if $M$ is simply connected, and a closed
$1$-form is exact if and only if all its periods vanish.
\end{proposition}

\begin{proof}
Compose the isomorphism $H^1_{\mathrm{dR}}(M) \cong H^1(M; \RR)$ of Theorem~\ref{de-rham:theorem-de-rham}, given by integration over smooth simplices, with $H^1(M; \RR) \cong \Hom_\RR(H_1(M; \RR), \RR) = \Hom(H_1(M), \RR)$
(Theorem~\ref{singular-cohomology:theorem-universal-coefficients} over the field $\RR$) and with $\Hom(H_1(M), \RR) \cong \Hom(\pi_1(M, x_0), \RR)$, which holds by the Hurewicz theorem
(Theorem~\ref{singular-homology:theorem-hurewicz}) because homomorphisms to the abelian group $\RR$ factor through the abelianisation. The composite is the period map.

The last statement can also be seen directly. If $\omega = df$, then $\int_\gamma\omega = f(\gamma(1)) - f(\gamma(0)) = 0$ for loops. Conversely, if all periods vanish, $f(x) = \int_{x_0}^x\omega$ along any piecewise smooth path is
well defined, since two paths differ by a loop, and in a chart around $x$ that is a ball, integrating along segments shows $df = \omega$, by the Poincar\'e lemma on the ball
(Theorem~\ref{de-rham:theorem-homotopy-invariance}): there $\omega = dg$, and $f - g$ is constant.
\end{proof}

\begin{example}\label{de-rham:example-angle-form}
On $M = \RR^2 \setminus \{0\}$ the \emph{angle form} $\omega = \frac{x\,dy - y\,dx}{x^2 + y^2}$ is closed, and its period on the unit circle is $2\pi$, so it is not exact. Since $\pi_1(M) \cong \ZZ$, Proposition~\ref{de-rham:proposition-periods}
gives $H^1_{\mathrm{dR}}(M) = \RR[\omega]$: a closed $1$-form $\eta$ on $M$ is exact if and only if $\int_{S^1}\eta = 0$, and in general $\eta - \frac{1}{2\pi}\big(\int_{S^1}\eta\big)\omega$ is exact. On the simply connected
$\RR^3 \setminus \{0\}$, by contrast, every closed $1$-form is exact.
\end{example}

\begin{counterexample}\label{de-rham:counterexample-infinite-dimensional}
De Rham cohomology of a noncompact manifold can be infinite-dimensional. On $M = \CC \setminus \ZZ$ let $\alpha_n$ be the angle form centred at $n$, the pullback of the angle form of
Counterexample~\ref{differential-forms:counterexample-closed-not-exact} by $z \mapsto z - n$. It is closed, and its integral over the circle $\gamma_m$ of radius $\frac12$ around $m$ is $2\pi$ if $m = n$ and $0$ otherwise (for $m \neq n$ the circle lies in
a half-plane, where $\alpha_n$ is exact). If $\sum_nc_n\alpha_n = d f$ for finitely many nonzero $c_n$, integrating over $\gamma_m$ gives $2\pi c_m = 0$ (Proposition~\ref{differential-forms:proposition-line-integrals}). So the classes
$[\alpha_n]$, $n \in \ZZ$, are linearly independent in $H^1_{\mathrm{dR}}(M)$.
\end{counterexample}

\begin{proposition}\label{de-rham:proposition-finite-quotient}
Let a finite group $G$ act freely on $M$ by diffeomorphisms, and $\pi \colon M \to M/G$ the quotient (Theorem~\ref{smooth-manifolds:theorem-smooth-quotient}). Then $\pi^*$ induces an isomorphism
$H^k_{\mathrm{dR}}(M/G) \cong H^k_{\mathrm{dR}}(M)^G$ onto the classes fixed by $G$.
\end{proposition}

\begin{proof}
As $\pi$ is a local diffeomorphism and $\pi \circ g = \pi$, $\pi^*$ identifies $\Omega^\bullet(M/G)$ with the complex $\Omega^\bullet(M)^G$ of $G$-invariant forms: an invariant form descends, because locally $\pi$ has smooth inverses that differ
by elements of $G$. The averaging operator $A = \frac{1}{|G|}\sum_{g \in G}g^*$ commutes with $d$ and is a projection of $\Omega^\bullet(M)$ onto $\Omega^\bullet(M)^G$. The map $H^k(\Omega^\bullet(M)^G) \to H^k_{\mathrm{dR}}(M)^G$ is injective: if an
invariant $\omega$ equals $d\eta$, then $\omega = A\omega = d(A\eta)$ with $A\eta$ invariant. It is surjective: if $[\omega]$ is $G$-fixed, then $A\omega$ is invariant, closed, and cohomologous to $\omega$, since each $g^*\omega$ is.
\end{proof}

\begin{example}\label{de-rham:example-real-projective}
$H^k_{\mathrm{dR}}(\RR\PP^n) \cong \RR$ for $k = 0$, and for $k = n$ if $n$ is odd; it vanishes otherwise. By Proposition~\ref{de-rham:proposition-finite-quotient} and Example~\ref{de-rham:example-spheres}, only the action of the
antipodal map $a$ on $H^n_{\mathrm{dR}}(S^n) \cong \RR$ matters. The volume form of $S^n$ is $\iota_X\Omega$ restricted to $S^n$, where $\Omega = dx^1 \wedge \cdots \wedge dx^{n+1}$ and $X = \sum x^i\partial_i$; as $a$ is the restriction of $-\mathrm{id}$,
$a^*\Omega = (-1)^{n+1}\Omega$ and $a$ preserves $X$, so $a^*$ acts on $H^n$ by $(-1)^{n+1}$, and the invariant part is $H^n$ for $n$ odd and $0$ for $n$ even. In particular $\RR\PP^n$ is not orientable for $n$ even, by
Corollary~\ref{de-rham:corollary-consequences}.
\end{example}

\begin{theorem}[K\"unneth formula]\label{de-rham:theorem-kunneth}
If $M$ has a finite good cover, the map $\bigoplus_{p+q=k}H^p_{\mathrm{dR}}(M) \otimes H^q_{\mathrm{dR}}(N) \to H^k_{\mathrm{dR}}(M \times N)$, $[\alpha] \otimes [\beta] \mapsto [\mathrm{pr}_M^*\alpha \wedge \mathrm{pr}_N^*\beta]$, is an isomorphism for every
manifold $N$.
\end{theorem}

\begin{proof}
Write $\psi_M$ for the map. It is natural in $M$ for restrictions to open subsets. If $M \cong \RR^m$, then $\mathrm{pr}_N \colon M \times N \to N$ is a homotopy equivalence and $H^\bullet_{\mathrm{dR}}(M) = \RR$ in degree $0$, so $\psi_M$ is
$\mathrm{pr}_N^*$, an isomorphism by Theorem~\ref{de-rham:theorem-homotopy-invariance}. For open $U, V \subseteq M$, tensoring the Mayer--Vietoris sequence of $U, V$ with the vector space $H^q_{\mathrm{dR}}(N)$ and summing over $q$ gives an
exact sequence, which $\psi$ maps to the Mayer--Vietoris sequence of $U \times N$, $V \times N$. The squares involving restrictions commute, and so does the one with the connecting maps: with a partition of unity $\{\rho_U, \rho_V\}$ on
$U \cup V$, the connecting maps are $[\omega] \mapsto [d(\rho_V\omega)]$ on both sides, and for closed $\beta$, $d(\mathrm{pr}_M^*(\rho_V\omega) \wedge \mathrm{pr}_N^*\beta) = \mathrm{pr}_M^*d(\rho_V\omega) \wedge \mathrm{pr}_N^*\beta$. So if $\psi$ is an isomorphism for $U$, $V$ and $U \cap V$, it is one for $U \cup V$
by the five lemma (Lemma~\ref{homological-algebra:lemma-five}). Now let $U_1, \ldots, U_k$ be a finite good cover of $M$; we show by induction on $k$ that $\psi$ is an isomorphism for every open set covered by $k$ open sets forming a good
cover. For $k = 1$ this is the case $\RR^m$. For the induction step apply the five lemma to $W = U_1 \cup \cdots \cup U_{k-1}$ and $U_k$, noting that $W \cap U_k$ is covered by the $k - 1$ sets $U_i \cap U_k$, which form a good cover.
\end{proof}

\begin{corollary}\label{de-rham:corollary-torus}
$H^\bullet_{\mathrm{dR}}(T^n)$ is the exterior algebra on the classes of $dx^1, \ldots, dx^n$, so $b_k(T^n) = \binom nk$; and $H^\bullet_{\mathrm{dR}}(S^1 \times S^2)$ has dimensions $1, 1, 1, 1$.
\end{corollary}

\begin{proof}
By Example~\ref{de-rham:example-spheres}, $H^\bullet(S^1)$ is spanned by $1$ and $[d\theta]$, and $T^n = (S^1)^n$ has a finite good cover; Theorem~\ref{de-rham:theorem-kunneth} and induction give the tensor product of $n$ copies of
$\RR \oplus \RR[d\theta]$, which is the exterior algebra because $d\theta \wedge d\theta = 0$ and the wedge product is graded commutative. The second statement is the same computation with $H^\bullet(S^2)$.
\end{proof}

\begin{theorem}[Poincar\'e duality in de Rham cohomology]\label{de-rham:theorem-poincare-duality-de-rham}
Let $M$ be an oriented $n$-manifold. The pairing
\[
H^k_{\mathrm{dR}}(M) \times H^{n-k}_{\mathrm{dR},c}(M) \to \RR, \qquad ([\alpha], [\beta]) \mapsto \int_M\alpha \wedge \beta,
\]
is well defined and nondegenerate if $M$ has a finite good cover (for instance if $M$ is compact); in that case $H^k_{\mathrm{dR}}(M) \cong H^{n-k}_{\mathrm{dR},c}(M)^*$. For compact oriented $M$ this is the
statement that $\int_M\alpha \wedge \beta$ is a perfect pairing between $H^k_{\mathrm{dR}}(M)$ and $H^{n-k}_{\mathrm{dR}}(M)$.
\end{theorem}

\begin{proof}
The pairing is well defined by Stokes' theorem: if $\alpha$ is closed and $\beta = d\gamma$ has compact support, then $\alpha \wedge \beta = \pm d(\alpha \wedge \gamma)$ has compact support and integrates to zero.
Nondegeneracy is proved by induction on the number of sets in a finite good cover, one in which all nonempty intersections are diffeomorphic to $\RR^n$ (such covers exist on any manifold, and finite ones
on compact manifolds), comparing the Mayer--Vietoris sequences of $H^\bullet_{\mathrm{dR}}$ and $H^\bullet_{\mathrm{dR},c}$ with the five lemma; see \cite[Chapter I, Section 5]{bott-tu}. Alternatively it follows from Theorem~\ref{de-rham:theorem-de-rham} and topological Poincar\'e duality
(Corollary~\ref{poincare-duality:corollary-cup-pairing}), since the wedge product corresponds to the cup product.
\end{proof}

\begin{theorem}[Cartan--Chevalley--Eilenberg]\label{de-rham:theorem-lie-group-cohomology}
Let $G$ be a compact connected Lie group. Every de Rham class of $G$ has a unique representative that is invariant under left and right translations, and the bi-invariant forms are closed. Consequently
$H^\bullet_{\mathrm{dR}}(G) \cong (\Lambda^\bullet\mathfrak{g}^*)^{G}$, the space of $\Ad(G)$-invariant alternating forms on the Lie algebra.
\end{theorem}

\begin{proof}
$G \times G$ acts on $G$ by $(a, b) \cdot x = axb^{-1}$; let $\mu$ be the Haar probability measure on the compact connected group $K = G \times G$ (Theorem~\ref{lie-groups:theorem-haar}) and, for a form $\omega$, let
$A\omega = \int_Kk^*\omega\,d\mu(k)$, defined pointwise; $A\omega$ is smooth and bi-invariant, $A$ commutes with $d$ (differentiation under the integral sign), and $A\omega = \omega$ for bi-invariant $\omega$.

\emph{Averaging does not change classes.} Let $\omega$ be closed. Each $k \in K$ is connected to the identity by a smooth path, so $k^*\omega$ is cohomologous to $\omega$ (Theorem~\ref{de-rham:theorem-homotopy-invariance}). For a closed form $\eta$
of complementary degree, Fubini's theorem gives $\int_GA\omega \wedge \eta = \int_K(\int_Gk^*\omega \wedge \eta)\,d\mu(k) = \int_G\omega \wedge \eta$, since the inner integral depends only on the class of $k^*\omega$ (Stokes). By Poincar\'e duality
(Theorem~\ref{de-rham:theorem-poincare-duality-de-rham}), $[A\omega] = [\omega]$.

\emph{Bi-invariant forms are closed.} The inversion $\iota(x) = x^{-1}$ satisfies $\iota \circ L_a = R_{a^{-1}} \circ \iota$, so it maps bi-invariant forms to bi-invariant forms, and $d_e\iota = -\mathrm{id}$ (differentiate $\exp(tX)^{-1} = \exp(-tX)$).
A bi-invariant form is determined by its value at $e$, so $\iota^*\omega = (-1)^k\omega$ for a bi-invariant $k$-form. As $d\omega$ is bi-invariant of degree $k + 1$, $(-1)^{k+1}d\omega = \iota^*d\omega = d\iota^*\omega = (-1)^kd\omega$, so $d\omega = 0$.

\emph{Uniqueness.} If a bi-invariant $\omega$ is exact, $\omega = d\eta$, then $\omega = A\omega = d(A\eta) = 0$, because $A\eta$ is bi-invariant, hence closed. So bi-invariant forms map isomorphically onto $H^\bullet_{\mathrm{dR}}(G)$. A
left-invariant form is determined by its value $\omega_e \in \Lambda^\bullet\mathfrak{g}^*$, and it is also right-invariant if and only if $\omega_e$ is invariant under $\Ad(g) = d_e(L_g \circ R_{g^{-1}})$.
\end{proof}

\begin{example}\label{de-rham:example-su2}
For $G = \mathrm{SU}(2)$, $\mathfrak{g} = \mathfrak{su}(2) \cong \RR^3$ and $\Ad$ maps $\mathrm{SU}(2)$ onto $\mathrm{SO}(3)$ (Exercise~\ref{lie-groups:exercise-su2}). The $\mathrm{SO}(3)$-invariants in $\Lambda^\bullet\RR^3$ are the constants and the
multiples of the determinant $e_1^* \wedge e_2^* \wedge e_3^*$, since $\Lambda^1$ and $\Lambda^2 \cong \RR^3$ are irreducible nontrivial representations. So $H^\bullet(\mathrm{SU}(2))$ is $\RR$ in degrees $0$ and $3$, in accordance with
$\mathrm{SU}(2) \cong S^3$.
\end{example}

\section{Currents}\label{de-rham:section-currents}

\begin{definition}\label{de-rham:definition-current}
Let $M$ be an oriented $n$-manifold. A \emph{$k$-current} is a linear functional $T$ on $\Omega^{n-k}_c(M)$ that is continuous in the sense of distributions: $T(\varphi_j) \to 0$ whenever the coefficients of
$\varphi_j$, in any chart, have supports in a fixed compact set and converge to $0$ together with all derivatives. Currents form a space $\mathcal{D}'_k(M)$, and $dT$ is defined by
$(dT)(\varphi) = (-1)^{k+1}T(d\varphi)$, so that $d \circ d = 0$.
\end{definition}

\begin{example}\label{de-rham:example-currents}
\begin{enumerate}
\item A locally integrable $k$-form $\eta$ defines a current $T_\eta(\varphi) = \int_M\eta \wedge \varphi$, and $dT_\eta = T_{d\eta}$ for smooth $\eta$ by Stokes' theorem; so currents generalise forms in the
same way as distributions generalise functions (Chapter~\ref{distributions}).
\item A compact oriented $k$-dimensional submanifold $Z \subseteq M$ defines the \emph{current of integration} $[Z](\varphi) = \int_Z\varphi$ for $\varphi \in \Omega^k_c(M)$. Stokes' theorem gives $d[Z] = 0$ when
$\partial Z = \emptyset$, so $[Z]$ has a cohomology class.
\end{enumerate}
\end{example}

\begin{theorem}\label{de-rham:theorem-currents}
On an oriented manifold the complex of currents computes de Rham cohomology: the inclusion of smooth forms into currents is a quasi-isomorphism. Under this identification and Poincar\'e duality, the class
of $[Z]$ for a compact oriented $k$-dimensional submanifold $Z$ of a compact oriented $n$-manifold $M$ is the Poincar\'e dual of $Z$: it is the unique class $\eta_Z \in H^{n-k}_{\mathrm{dR}}(M)$ with
\[
\int_Z\varphi = \int_M\eta_Z \wedge \varphi \qquad \text{for all closed } \varphi \in \Omega^k(M).
\]
\end{theorem}

This is proved in \cite[Chapter IV]{de-rham-book} and \cite[Chapter 0, Section 4]{griffiths-harris}; the first statement is a sheaf-theoretic argument like the proof of
Theorem~\ref{de-rham:theorem-de-rham}, applied to the fine resolution of $\RR_M$ by sheaves of currents, and the second follows by comparing both sides with the topological Poincar\'e dual
(Remark~\ref{poincare-duality:remark-fundamental-classes-subspaces}).

\begin{remark}\label{de-rham:remark-cycle-classes}
Currents are the natural home for the cycle classes of Hodge theory: an analytic subvariety $Z$ of a compact complex manifold, possibly singular, defines a current of integration because its singular locus
has measure zero and $\int_{Z_{\mathrm{reg}}}\varphi$ converges (Chapter~\ref{cycle-class}); its class in $H^{2n-2k}(X; \ZZ)$ is the cycle class whose Hodge-theoretic characterisation the Hodge conjecture
concerns.
\end{remark}

\section{Exercises}\label{de-rham:section-exercises}

\begin{exercise}\label{de-rham:exercise-torus}
Compute $H^\bullet_{\mathrm{dR}}(T^n)$ and show that it is the exterior algebra on $H^1$, generated by the forms $d\theta^1, \ldots, d\theta^n$.
\end{exercise}

\begin{solution}
By Theorem~\ref{de-rham:theorem-de-rham} and Example~\ref{singular-cohomology:example-torus}, $H^\bullet_{\mathrm{dR}}(T^n) \cong \Lambda^\bullet\RR^n$. The classes of the closed forms $d\theta^i$, which are
not exact since $\int_{S^1_i}d\theta^i = 2\pi$, generate: their products are the $\binom{n}{k}$ linearly independent classes of $d\theta^{i_1} \wedge \cdots \wedge d\theta^{i_k}$, which can be detected by
integration over the corresponding subtori.
\end{solution}

\begin{exercise}\label{de-rham:exercise-compactly-supported-line}
Show that $H^0_{\mathrm{dR},c}(\RR) = 0$ and $H^1_{\mathrm{dR},c}(\RR) \cong \RR$, with the isomorphism given by integration; so compactly supported de Rham cohomology is not homotopy invariant.
\end{exercise}

\begin{solution}
A compactly supported function with $df = 0$ is constant, hence zero. The map $\int_\RR \colon \Omega^1_c(\RR) \to \RR$ is surjective, and $f\,dx$ integrates to zero if and only if $F(x) = \int_{-\infty}^xf$ has
compact support, in which case $f\,dx = dF$ is exact.
\end{solution}

\begin{exercise}\label{de-rham:exercise-degree-de-rham}
Let $f \colon M \to N$ be a smooth map of compact connected oriented $n$-manifolds. Show that $f^* \colon H^n_{\mathrm{dR}}(N) \to H^n_{\mathrm{dR}}(M)$ is multiplication by the degree of $f$, and that a map of
nonzero degree is surjective.
\end{exercise}

\begin{exercise}\label{de-rham:exercise-kunneth-de-rham}
Prove the K\"unneth formula $H^\bullet_{\mathrm{dR}}(M \times N) \cong H^\bullet_{\mathrm{dR}}(M) \otimes H^\bullet_{\mathrm{dR}}(N)$ for manifolds with finite good covers, by induction with Mayer--Vietoris.
\end{exercise}

\begin{exercise}\label{de-rham:exercise-punctured-plane}
Let $M = \RR^2 \setminus \{p_1, \ldots, p_k\}$ for distinct points $p_i$. Show that $H^1_{\mathrm{dR}}(M) \cong \RR^k$, with a basis given by the angle forms $\alpha_i$ centred at the $p_i$,
and that a closed $1$-form $\eta$ on $M$ is exact if and only if its integrals over small circles around all the $p_i$ vanish.
\end{exercise}

\begin{solution}
After a rotation we may assume the $p_i$ have distinct first coordinates, $p_1$ having the largest. Let $U$ be $M$ intersected with $\{x < a\}$ and $V$ be $M$ intersected with
$\{x > b\}$, where $b < a$ are chosen so that $p_1$ is the only point with first coordinate $> b$ and none has first coordinate in $[b, a]$. Then $U \cap V$ is a strip, which is
contractible, so Mayer--Vietoris gives $H^1(M) \cong H^1(U) \oplus H^1(V)$. The half-planes are diffeomorphic to $\RR^2$, so $V$ is a plane minus one point, with $H^1 \cong \RR$ spanned by
$\alpha_1$ (Example~\ref{de-rham:example-angle-form}), and $U$ is a plane minus $k - 1$ points; induction gives $\RR^k$. The integral of $\alpha_i$ over a small circle around $p_j$ is
$2\pi\delta_{ij}$ (away from $p_i$, $\alpha_i$ is exact on a disc), so the map $[\eta] \mapsto (\int_{C_j}\eta)_j$ is surjective onto $\RR^k$, hence an isomorphism, and the $[\alpha_i]$ form a basis.
\end{solution}

\begin{exercise}\label{de-rham:exercise-solid-angle}
On $\RR^3 \setminus \{0\}$ let $\sigma = (x\,dy \wedge dz + y\,dz \wedge dx + z\,dx \wedge dy)/r^3$, $r = (x^2 + y^2 + z^2)^{1/2}$ (the \emph{solid angle form}). Show that $\sigma$ is closed, that
$\int_{S^2}\sigma = 4\pi$, and deduce that $H^2_{\mathrm{dR}}(\RR^3 \setminus \{0\}) = \RR[\sigma]$ and that $\sigma$ is not exact.
\end{exercise}

\begin{solution}
By Example~\ref{differential-forms:example-forms-r3}, $d\sigma = \operatorname{div}(x/r^3, y/r^3, z/r^3)\,dx \wedge dy \wedge dz$, and $\partial_x(x/r^3) = 1/r^3 - 3x^2/r^5$, so the divergence is
$3/r^3 - 3r^2/r^5 = 0$. On the unit sphere $r = 1$ and $\sigma$ restricts to the area form (Example~\ref{differential-forms:example-integrals}(2)), so $\int_{S^2}\sigma = 4\pi$. An exact form would have
integral $0$ over the closed surface $S^2$ (Theorem~\ref{differential-forms:theorem-stokes}). Finally $\RR^3 \setminus \{0\}$ is homotopy equivalent to $S^2$, so its $H^2$ is one-dimensional
(Example~\ref{de-rham:example-spheres}), spanned by the nonzero class $[\sigma]$.
\end{solution}

\begin{exercise}\label{de-rham:exercise-explicit-primitive}
Decide whether the $2$-form $\omega = z\,dx \wedge dy + x\,dy \wedge dz + y\,dz \wedge dx$ on $\RR^3$ is closed, and find a primitive of the closed form
$\omega^{\prime} = dx \wedge dy + dy \wedge dz$ using Corollary~\ref{de-rham:corollary-explicit-poincare}.
\end{exercise}

\begin{solution}
$d\omega = \partial_zz\,dz \wedge dx \wedge dy + \partial_xx\,dx \wedge dy \wedge dz + \partial_yy\,dy \wedge dz \wedge dx = 3\,dx \wedge dy \wedge dz \neq 0$, so $\omega$ is not closed and has no primitive.
For $\omega^{\prime}$, which has constant coefficients and is closed, $\iota_X\omega^{\prime} = x\,dy - y\,dx + y\,dz - z\,dy$ and $\eta = \frac12(x\,dy - y\,dx + y\,dz - z\,dy)$; then
$d\eta = \frac12(2\,dx \wedge dy + 2\,dy \wedge dz) = \omega^{\prime}$.
\end{solution}
